\documentclass[12pt,a4paper]{article}
\usepackage[T2A]{fontenc}
\usepackage[utf8]{inputenc}
\usepackage[russian]{babel}
\usepackage{geometry}
\usepackage{indentfirst}
\usepackage{enumitem}

\geometry{left=25mm,right=15mm,top=20mm,bottom=20mm}

\setlength{\parindent}{1.25cm}
\setlength{\parskip}{0pt}

\pagestyle{plain}

\begin{document}

\begin{center}
{\Large \textbf{Функциональная спецификация}}\\[0.4cm]
{\large курсовой работы}\\[0.6cm]
{\large Программа автоматической синхронизации файлов в двух каталогах на разных компьютерах (клиентская синхронизация)}
\end{center}

\vspace{0.5cm}

\section{Назначение программы}

Разрабатываемая программа представляет собой приложение для автоматической синхронизации файлов между двумя каталогами, расположенными на разных компьютерах. Программа предназначена для пользователей, которым необходимо поддерживать одинаковое содержимое каталогов на двух компьютерах.

Работа с программой осуществляется через графический интерфейс, реализованный с использованием библиотеки Java Swing. На каждом компьютере запускается экземпляр программы, который взаимодействует с другим экземпляром по сети посредством TCP-сокетов.

\section{Запуск и настройка программы}

После запуска приложения пользователю отображается главное окно программы.

Перед началом работы пользователь задаёт следующие параметры:

\begin{itemize}[nosep]
    \item путь к каталогу, который необходимо синхронизировать;
    \item IP-адрес партнёра;
    \item номер порта партнёра;
    \item локальный порт для входящих соединений.
\end{itemize}

После ввода параметров пользователь нажимает кнопку \textbf{«Подключиться»}. При успешном подключении программа готова к выполнению синхронизации.

\section{Синхронизация файлов}

При выполнении синхронизации экземпляры программы соединяются по TCP-сокету и обмениваются списками файлов своих каталогов. Для каждого файла передаются относительный путь, размер и хэш SHA-256.

По полученным спискам каждый экземпляр определяет файлы, отсутствующие в его каталоге, и копирует их от партнёра. По завершении синхронизации оба каталога содержат одинаковый набор файлов.

Для синхронизации реализуются следующие операции:

\begin{itemize}[nosep]
    \item установление соединения с партнёром;
    \item обмен списками файлов каталогов;
    \item вычисление хэшей файлов (SHA-256);
    \item копирование недостающих файлов;
    \item завершение соединения.
\end{itemize}

\section{Обнаружение конфликтов}

Если в двух каталогах присутствуют файлы с одинаковым относительным именем, но различным содержимым (различаются хэши), такая ситуация считается конфликтом.

При обнаружении конфликта файл не перезаписывается. В журнал работы заносится сообщение с указанием имени файла и хэшей обеих версий.

\section{Журнал работы}

Все действия программы записываются в журнал работы (\texttt{sync.log}):

\begin{itemize}[nosep]
    \item подключение к партнёру;
    \item начало и завершение синхронизации;
    \item список скопированных файлов;
    \item обнаруженные конфликты;
    \item ошибки соединения.
\end{itemize}

Основные сообщения одновременно отображаются в окне программы.

\section{Хранение данных}

Информация о состоянии каталога и журнал работы программы хранятся в текстовых файлах. При выполнении операций синхронизации соответствующие данные обновляются. При повторном запуске программы сохранённые данные автоматически загружаются.

\section{Интерфейс программы}

Пользователь взаимодействует с программой посредством графического интерфейса, разработанного с использованием библиотеки Java Swing.

Главное окно программы содержит:

\begin{itemize}[nosep]
    \item поле выбора каталога синхронизации;
    \item кнопку выбора каталога (\texttt{JFileChooser});
    \item поле ввода IP-адреса партнёра;
    \item поле ввода порта партнёра;
    \item поле ввода локального порта;
    \item кнопку \textbf{«Подключиться»};
    \item кнопку \textbf{«Синхронизировать»};
    \item область отображения журнала работы;
    \item строку состояния.
\end{itemize}

Во время выполнения синхронизации пользователь может наблюдать процесс подключения, обмена данными и копирования файлов. После завершения операции отображается итоговое сообщение с количеством переданных файлов и обнаруженных конфликтов.

При возникновении ошибок выводятся диалоговые окна с описанием причины ошибки.

\subsection*{Главное окно программы}

Главное окно обеспечивает выполнение всех основных операций программы без использования командной строки. Пользователь выбирает каталог, вводит параметры подключения и запускает синхронизацию нажатием соответствующей кнопки. Журнал работы отображается непосредственно в окне приложения.

\section{Функции базовой версии (оценка 4)}

В базовой версии программы пользователю доступны следующие функции:

\begin{itemize}[nosep]
    \item запуск программы через графический интерфейс Swing;
    \item выбор каталога синхронизации посредством диалога выбора файлов;
    \item ввод сетевых параметров подключения;
    \item формирование снимка состояния каталога (размер и хэш SHA-256 каждого файла);
    \item запуск синхронизации нажатием кнопки;
    \item обмен списками файлов между экземплярами программы;
    \item копирование отсутствующих файлов;
    \item обнаружение конфликтов и запись их в журнал;
    \item отображение журнала работы в окне программы;
    \item вывод итоговой информации о выполненной синхронизации;
    \item обработка ошибок с отображением понятных сообщений пользователю.
\end{itemize}

\section{Основные сценарии работы}

Основной сценарий работы пользователя:

\begin{enumerate}[nosep]
    \item Пользователь запускает программу на обоих компьютерах.
    \item Выбирает каталог синхронизации.
    \item Вводит IP-адрес и номера портов.
    \item Нажимает кнопку \textbf{«Подключиться»}.
    \item После успешного подключения нажимает кнопку \textbf{«Синхронизировать»}.
    \item Экземпляры программы обмениваются списками файлов.
    \item Недостающие файлы копируются от партнёра.
    \item При обнаружении конфликта информация заносится в журнал.
    \item Пользователь получает итоговый отчёт о синхронизации.
    \item Пользователь завершает работу программы.
\end{enumerate}

Таким образом, программа позволяет автоматически синхронизировать каталоги между двумя компьютерами, контролировать конфликты файлов и отслеживать ход выполнения операций посредством графического интерфейса.

\section{Требования к реализации и отчётности}

Программа реализуется на языке Java (JDK 17 и выше) с использованием только стандартной библиотеки.

Архитектура приложения построена по шаблону Model--View--Controller. Проект разделён на пакеты:

\begin{itemize}[nosep]
    \item \texttt{model};
    \item \texttt{view} (графический интерфейс Swing);
    \item \texttt{controller};
    \item \texttt{network};
    \item \texttt{util}.
\end{itemize}

Графический интерфейс реализован средствами Java Swing и отделён от бизнес-логики приложения в соответствии с архитектурой MVC.

Сборка осуществляется с помощью Maven. Выполняются компиляция проекта, запуск unit-тестов модели и упаковка приложения в исполняемый JAR-файл.

Проект включает unit-тесты (JUnit~5), проверяющие построение снимка каталога, вычисление хэшей файлов и построение плана синхронизации.

Отчёт оформляется в \LaTeX{} и включает титульный лист, функциональную спецификацию, руководство пользователя со скриншотами графического интерфейса, описание архитектуры программы и описание основных классов.

\end{document}